\section*{September 4, 2026: A common distance comparison retains only six treated sites}
\textit{Agent-drafted follow-up to the distance-ring diagnostic.}

Jacob requested a comparison holding the school roster and overlap-exclusion
radius fixed across distances. The half-mile exclusion is now applied in
every band: each home must be more than half a mile from any second study
site, welcoming site, or other candidate site. The comparison starts with the
selected 30 closed and 49 still-open schools
and the 2008--2018 housing-price sample. Within those eligible homes, the
three bands are 0--0.125, 0.125--0.25, and 0.25--0.5 miles from the nearest
study site.

To prevent empty inner rings from changing the annual comparison, a
site-year enters only when it has at least one sale in every band. A site
must have at least one such year before 2013 and one after 2013. This gives
the same six treated and 26 control sites in all bands, with the same
contributing sites across bands within each year. The yearly roster still
varies: three to six treated sites and 21 to 25 control sites contribute.
There are 52 treated and 252 control site-years in each band. No missing
prices are imputed. Requiring complete coverage in every year would leave
only three treated sites; that additional restriction is not imposed.

The retained treated/control transaction counts are 211/1,079 in the closest
band, 398/2,427 in the middle band, and 815/7,108 in the outer band. The strict
half-mile overlap rule alone leaves only seven treated sites with any sales
in the closest band. Most of the loss therefore precedes the requirement
for common annual coverage. This diagnostic describes a small subset of the
original closure comparison, not the effect for all 30 closed schools.

\begin{center}
\scriptsize
\input{../output/common_ring_table.tex}
\end{center}
{\footnotesize Annual DiDs with housing controls and site/year fixed effects;
equal observed site-year weights. Pre: 2008--2012; post: 2014--2018; 2013
omitted. Brackets are 95\% site-clustered intervals. Prices use 2022 dollars;
log coefficients and endpoints are converted using $100[\exp(\beta)-1]$.}

The point estimates do not reveal a clear effect concentrated closest to
the schools. Dollar effects are about \$31,000, \$38,000, and \$23,000;
log effects imply approximately 6\%, 8\%, and 16\%. Every interval includes
zero and economically large effects. The raw treated means jump in 2015
across all three bands, but only three treated sites contribute that year.
Matching the annual roster across bands does not eliminate changes in that
roster over time or changes in the homes sold within each site.

Replacing site intercepts with a treated/control group intercept remains
consequential. In the middle band, the site-weighted dollar estimate changes
from about \$38,000 to minus \$4,000. Both intervals are wide. School-site
fixed effects address persistent differences among the retained sites;
they do not resolve differential neighborhood trends or selection into sale.
The small number of treated sites makes the conventional clustered
intervals and pretrend tests weak evidence for choosing a specification.

This exercise removes across-band changes in the site roster and exclusion
rule. Differences across bands cannot be read as how the causal effect varies
with distance: schools qualify partly on the basis of sales after closure,
different schools contribute in different years, and prices are observed
only for homes that sell. Codex's assessment is to retain this as an audit,
rather than adopt the severely restricted sample as the main design.

\clearpage
\begin{center}
\includegraphics[width=\textwidth,height=.93\textheight,keepaspectratio]{../input/ring_common_raw_equal_site_year.png}
\end{center}
\clearpage
\begin{center}
\includegraphics[width=\textwidth,height=.93\textheight,keepaspectratio]{../input/ring_common_event_hedonics.png}
\end{center}
\clearpage
\begin{center}
\includegraphics[width=\textwidth,height=.93\textheight,keepaspectratio]{../input/ring_common_event_fixed_effects.png}
\end{center}
\clearpage
\begin{center}
\includegraphics[width=\textwidth]{../input/ring_common_did_by_ring.png}
\end{center}
{\footnotesize Reproduction: Make in the home-price-rings audit code directory,
then the logbook code directory. The common-design outputs retain the same
120-model family as the earlier audit; cumulative buffers combine the same
eligible site-years. Revision c200c31 plus working changes; existing cleaned
2008--2018 sales and school distances. Independent cell tables reproduce all
five sample memberships. Dummy-variable weighted least squares reproduces
controlled DiDs within \$0.000001; unadjusted site-weighted event studies
match regressions on site-year means. Earlier ring CSV outputs remain
byte-for-byte unchanged. The saved support grid records every excluded or
empty site-band-year cell.}
